Python 编程语言拥有独特的优势与魅力，但往往难以捉摸。许多熟悉其他语言的程序员，在用一种有限的心态接触 Python，而不是去接纳它的全部能力。也有一些程序员走向了另一个极端，过度使用某些 Python 特性，从而为日后埋下重大问题。

本书深入阐释了编写程序的 Pythonic（地道写法）之道：使用 Python 的最佳方式。它建立在一种我假定你已经具备的、对该语言的基本理解之上。初级程序员将学到 Python 关键特性方面的最佳实践；有经验的程序员则将学会如何满怀信心地去接纳一门新工具。

通过本书，我希望能帮助你用 Python 去达成你的目标——无论那目标是什么；或者至少，帮助你在编程之旅中收获更多乐趣。

本书每一章都包含一组宽泛但彼此相关的条目（item）。你可以随意在条目之间跳转，跟随自己的兴趣阅读。每个条目都包含简洁而具体的指导，解释你如何才能更有效地编写 Python 程序。条目中既有该做什么的建议，也有不该做什么的提醒，还有如何把握恰当平衡、以及为何这是最佳选择的说明。条目之间会互相引用，以便你在阅读时更容易填补知识空白。

本第三版涵盖了截至 Python 3.13 版本的语言特性（参见 Item 1「了解你正在使用的 Python 版本」）。相比第二版，本书新增了 35 个全新条目。第二版中的大多数条目都经过修订并被保留，但其中许多都经历了大幅更新。对于某些条目，由于过去五年间随着 Python 日渐成熟、最佳实践不断演进，我的建议已经彻底改变。

Python 对其标准库（standard library）采取了「batteries included（开箱即用）」的理念。这些内置（built-in）包中有许多与 Pythonic 写法如此紧密地交织在一起，以至于它们几乎成了语言规范（specification）的一部分。标准模块的总数过于庞大，无法在本书中全部涵盖，但我已收录了那些我认为必须了解并使用的模块。

Python 还拥有一个活跃的、由社区构建的模块生态，它们以有价值的方式扩展了这门语言。尽管我在各个条目中提及了值得了解的重要包，但本书无意成为一本详尽的参考手册。同样地，尽管 Python 的包管理十分重要，我仍避免深入其中的细节，因为它正处于快速变化与演进之中。

第 1 章：Pythonic 思维

Python 社区已经习惯用 Pythonic 这个形容词来描述遵循某种特定风格的代码。Python 的惯用法（idiom）是随着使用这门语言的经验、以及与其他程序员的协作，随着时间的推移而逐渐形成的。本章涵盖了在 Python 中做最常见之事的最佳方式。

第 2 章：字符串与切片

Python 为字符串（string）与序列处理提供了内置的语法、方法和模块。这些能力如此基础，以至于你几乎会在每一个程序中看到它们。它们使 Python 成为一门用于解析文本、检视数据格式，以及与计算机所用的底层二进制表示进行交互的绝佳语言。

第 3 章：循环与迭代器

对顺序数据进行遍历处理，是程序中的一项关键需求。对于涉及内置数据类型、容器（container）类型和用户自定义类的最常见任务，Python 中的循环（loop）用起来自然而强大。Python 还支持迭代器（iterator），它带来了一种更具函数式风格的方法，能以显著的收益来处理任意的数据流。

第 4 章：字典

Python 内置的字典（dictionary）类型，是程序中用于记账的一种万能数据结构（data structure）。与简单的列表（list）相比，字典在添加与删除条目时能提供好得多的性能。Python 还拥有特殊的语法和相关的内置模块，使字典的能力超越了你在其他语言中从哈希表（hash table）所预期的范围。

第 5 章：函数

Python 中的函数（function）拥有多种多样的额外特性，能让程序员的工作更轻松。其中一些与其他编程语言中的能力相似，但许多是 Python 独有的。本章讲解如何使用函数来表达意图、促进复用并减少缺陷（bug）。

第 6 章：推导式与生成器

Python 拥有特殊语法，可快速遍历列表（list）、字典（dictionary）和集合（set）以生成派生的数据结构。它还允许函数以增量方式返回一个可迭代（iterable）值的流。本章讲解这些特性如何带来更佳的性能、更少的内存占用以及更强的可读性。

第 7 章：类与接口

Python 是一门面向对象（object-oriented）的语言。在 Python 中把事情做成，往往需要编写新的类（class），并定义它们如何通过各自的接口（interface）与层次结构（hierarchy）相互交互。本章讲解如何使用类，借助对象来表达预期的行为。

第 8 章：元类与属性

元类（metaclass）与动态属性（attribute）是强大的 Python 特性。然而，它们也让你能够实现极其怪异、出人意料的行为。本章讲解使用这些机制的常见惯用法（idiom），以确保你遵循最小惊讶原则（rule of least surprise）。

第 9 章：并发与并行

借助线程（thread）和异步协程（coroutine）等特性，Python 让编写并发（concurrency）程序变得轻松，这类程序看似在同一时刻做许多不同的事情。Python 也可以通过系统调用、子进程（subprocess）和特殊模块来执行并行（parallelism）工作。本章讲解如何在这些微妙不同的情境中最好地运用 Python。

第 10 章：健壮性

让程序在遇到意外情况时依然可靠，与让程序功能正确同样重要。Python 拥有内置的特性和模块，有助于强化你的程序，使其在多种多样的情境下都保持健壮（robust）。

第 11 章：性能

Python 包含多种能力，能让程序以相对较低的投入，取得令人惊讶的出色性能。借助这些特性，可以在保留 Python 高层特性所带来的生产力增益的同时，从宿主系统中榨取最大性能。

第 12 章：数据结构与算法

Python 包含许多标准数据结构（data structure）与算法（algorithm）的优化实现，能帮助你以最小的投入获得高性能。这门语言还提供了久经考验的数据类型与辅助函数，用于常见任务（例如处理货币与时间），让你能专注于程序的核心需求。

第 13 章：测试与调试

无论用何种语言编写，你都应当始终测试你的代码；但对于 Python 而言，测试尤为重要。Python 的动态（dynamic）特性会以独特的方式增加运行时错误（runtime error）的风险。所幸，它们也让编写测试和诊断出故障的程序变得更容易。本章涵盖 Python 内置的测试与调试工具。

第 14 章：协作

在 Python 程序上协作，要求你审慎地对待自己编写代码的方式。即便你独自工作，也会希望了解如何使用他人编写的模块（module）。本章涵盖让众人能够在 Python 程序上协同工作的标准工具与最佳实践。

本书中的 Python 代码片段以 \texttt{monospace font} 显示，并带有语法高亮。当代码行很长时，我会使用 \texttt{'} 来标示换行发生的地方。我会用省略号 \texttt{...} 截断部分片段，用以表示那些存在但并非表达要点所必需的代码区域。你需要下载完整的示例代码（下文说明了获取方式），才能让这些被截断的片段在你的计算机上正确运行。

为了让代码示例更好地契合书籍的版式，或为了突出最重要的部分，我在 Python 风格指南上做了一些艺术化处理。我还在代码示例中省略了内嵌文档，以缩减其体积。我强烈建议你不要在自己的项目中效仿这种做法；相反，你应当遵循风格指南（参见 Item 2「遵循 PEP 8 风格指南」），并编写文档（参见 Item 118「为每个函数、类和模块编写文档字符串」）。

本书中大多数代码片段都附带了运行该代码所对应的输出。当我说「输出（output）」时，我指的是控制台（console）或终端（terminal）输出：也就是在交互式解释器（interpreter）中运行 Python 程序时你所看到的内容。输出部分以 \texttt{monospace font} 显示，并且每一段之前都有一行 \texttt{>>>}（即 Python 的交互式提示符）。

其用意在于，你可以把这些代码片段键入 Python 解释器（shell），并复现出预期的输出。

最后，还有一些以 \texttt{monospace font} 显示的段落，它们前面没有 \texttt{>>>} 这一行。这些代表的是 Python 解释器之外的程序运行输出。这些示例常常以 \texttt{\$} 字符开头，用以表示我正在从一个类似 Bash 的命令行解释器中运行程序。如果你在 Windows 或其他类型的系统上运行这些命令，可能需要相应地调整程序名与参数。

获取代码与勘误

把本书中的许多示例当作没有穿插散文的整体程序来阅读，是很有用的。这也给了你一个亲自把玩代码、理解程序为何如所述般运行的机会。你可以在本书的网站 \texttt{https://effectivepython.com} 上找到本书所有代码片段的源代码。该网站还提供了如何报告错误的说明。预先感谢你就所发现的任何错误与我联系。

注册你的副本

请在 InformIT 网站上注册你的《Effective Python：125 Specific Ways to Write Better Python，第三版》副本，以便在其可用时便捷地获取更新和（或）勘误。要开始注册流程，请访问 \texttt{informit.com/register} 并登录或创建账户。输入产品 ISBN（9780138172183）并点击 \texttt{Submit}。如果你希望在新版本与更新发布时收到独家优惠通知，请勾选相应复选框以接收我们发来的电子邮件。
